\documentclass[11pt,a4paper]{article}
\usepackage[T1]{fontenc}
\usepackage[utf8]{inputenc}
\usepackage[french]{babel}
\usepackage{lmodern}
\usepackage{geometry}
\geometry{top=1.8cm,bottom=2cm,left=2cm,right=2cm}
\usepackage{graphicx}
\usepackage{xcolor}
\usepackage{booktabs,longtable,tabularx,array}
\usepackage{enumitem}
\usepackage{listings}
\usepackage{hyperref}
\usepackage[most]{tcolorbox}
\usepackage{fancyhdr}
\usepackage{titlesec}
\usepackage{amsmath}
\usepackage{microtype}
\usepackage{parskip}
\usepackage{lastpage}
\definecolor{STGSBlue}{HTML}{16324F}
\definecolor{STGSTeal}{HTML}{0B7A75}
\definecolor{STGSGray}{HTML}{F3F5F7}
\definecolor{STGSBorder}{HTML}{CAD4DE}
\definecolor{STGSRed}{HTML}{A23B3B}
\definecolor{STGSGreen}{HTML}{276749}
\definecolor{STGSOrange}{HTML}{A15C16}
\hypersetup{colorlinks=true,linkcolor=STGSBlue,urlcolor=STGSTeal,citecolor=STGSTeal,pdftitle={STGS documentation technique}}
\pagestyle{fancy}
\fancyhf{}
\fancyhead[L]{\textcolor{STGSBlue}{\textbf{STGS - Documentation d'exécution}}}
\fancyhead[R]{\small Projet C++20 / Linux}
\fancyfoot[L]{\small Révision documentaire : 21 août 2026}
\fancyfoot[R]{\small Page \thepage/\pageref{LastPage}}
\setlength{\headheight}{14pt}
\titleformat{\section}{\Large\bfseries\color{STGSBlue}}{\thesection}{0.7em}{}
\titleformat{\subsection}{\large\bfseries\color{STGSTeal}}{\thesubsection}{0.7em}{}
\titleformat{\subsubsection}{\normalsize\bfseries\color{STGSBlue}}{\thesubsubsection}{0.7em}{}
\setlist[itemize]{leftmargin=1.5em,itemsep=0.2em,topsep=0.3em}
\setlist[enumerate]{leftmargin=1.7em,itemsep=0.25em,topsep=0.3em}
\lstdefinestyle{terminal}{
  backgroundcolor=\color{STGSGray},
  basicstyle=\ttfamily\small,
  frame=single,
  rulecolor=\color{STGSBorder},
  breaklines=true,
  columns=fullflexible,
  keepspaces=true,
  showstringspaces=false,
  xleftmargin=0.4em,
  xrightmargin=0.4em,
  aboveskip=0.8em,
  belowskip=0.8em
}
\lstset{style=terminal}
\newcommand{\code}[1]{\texttt{\detokenize{#1}}}
\newcommand{\source}[2]{\textcolor{STGSTeal}{\textbf{Source :}} \code{#1}, lignes #2}
\newtcolorbox{callflow}[1]{
  enhanced,breakable,colback=STGSGray,colframe=STGSBlue,
  title={#1},fonttitle=\bfseries,left=1.2mm,right=1.2mm,top=1mm,bottom=1mm
}
\newtcolorbox{limitbox}[1]{
  enhanced,breakable,colback=white,colframe=STGSOrange,
  title={#1},fonttitle=\bfseries,left=1.2mm,right=1.2mm,top=1mm,bottom=1mm
}
\newtcolorbox{whybox}[1]{
  enhanced,breakable,colback=white,colframe=STGSTeal,
  title={#1},fonttitle=\bfseries,left=1.2mm,right=1.2mm,top=1mm,bottom=1mm
}
\newtcolorbox{warningbox}[1]{
  enhanced,breakable,colback=white,colframe=STGSRed,
  title={#1},fonttitle=\bfseries,left=1.2mm,right=1.2mm,top=1mm,bottom=1mm
}
\newcommand{\cover}[3]{
\begin{titlepage}
\centering
\vspace*{0.7cm}
\includegraphics[width=0.88\textwidth]{assets/Banniere.png}\\[1.2cm]
{\Huge\bfseries\color{STGSBlue} #1\par}
\vspace{0.35cm}
{\Large #2\par}
\vspace{0.7cm}
\begin{tcolorbox}[colback=STGSGray,colframe=STGSBorder,width=0.86\textwidth]
#3
\end{tcolorbox}
\vfill
{\large Space Telemetry Ground Station (STGS)\par}
{\small Documentation basée sur l'archive \code{cpp-space-telemetry-ground-station-showcase-modular.zip}.\par}
{\small Les numéros de lignes correspondent à cette révision et peuvent évoluer après modification du code.\par}
\end{titlepage}
\tableofcontents
\newpage
}
\begin{document}
\cover{Volume 2 - Commandes de la station sol}{Traçage précis de \texttt{stgs\_ground\_station}}{
Chaque commande est suivie depuis le parsing CLI jusqu'au résumé final. Les différences TCP, UDP, replay et auto-port sont explicitées, ainsi que les options de pipeline, santé, filtrage et sortie.
}

\section{Point d'entrée commun}
Toutes les commandes de station commencent par :
\begin{callflow}{Appels communs}
\ttfamily
main(argc, argv)
$\rightarrow$ parseArgs(argc, argv)
$\rightarrow$ run(Options)
$\rightarrow$ TerminalUi
$\rightarrow$ Logger + StationHealthMonitor + PosixSignalStopController
$\rightarrow$ GroundStationPipeline::run()
\end{callflow}

\source{apps/ground_station.cpp}{15--27} ; \source{apps/ground_station/GroundStationApp.cpp}{107--124}

\section{Commande TCP classique}
\begin{lstlisting}
./build/stgs_ground_station --tcp --port 9000 --decoder-threads 4 --output telemetry.csv
\end{lstlisting}

\subsection{Parsing exact}
\code{parseArgs()} sélectionne \code{Transport::Tcp}, valide le port dans [1,65535], fixe 4 workers et le chemin de sortie. Le format est \code{Auto} ; l'extension n'étant pas \code{.json}, \code{resolveOutputFormat()} choisira CSV.

\subsection{Chaîne d'appels}
\begin{callflow}{TCP sur le port 9000}
\ttfamily
main
$\rightarrow$ parseArgs
$\rightarrow$ run
$\rightarrow$ GroundStationPipeline ctor
$\rightarrow$ TelemetryOutput ctor
$\rightarrow$ pipeline.run
$\rightarrow$ startWorkers
$\rightarrow$ receiveNetwork
$\rightarrow$ NetworkServer::run
$\rightarrow$ NetworkServer::runTcp
$\rightarrow$ createBoundSocket
$\rightarrow$ bind + listen + poll
$\rightarrow$ accept
$\rightarrow$ recv
$\rightarrow$ StreamFrameExtractor::feed
$\rightarrow$ callback
$\rightarrow$ submit
$\rightarrow$ rawQueue
$\rightarrow$ decoderLoop/decodeFrame
$\rightarrow$ decodedQueue
$\rightarrow$ writerLoop
$\rightarrow$ santé + STGA + CSV
\end{callflow}

\subsection{Pourquoi chaque étape existe}
\begin{longtable}{p{0.23\textwidth}p{0.32\textwidth}p{0.37\textwidth}}
\toprule
Fonction & Action & Pourquoi \\
\midrule
\endhead
\code{createBoundSocket()} & socket IPv4, SO\_REUSEADDR, bind & Centraliser l'initialisation et garantir le RAII en cas d'erreur. \\
\code{runTcp()} & non-blocking + listen + poll & Supporter plusieurs clients et un arrêt coopératif. \\
\code{StreamFrameExtractor::feed()} & reconstruit les frontières & TCP ne préserve pas les messages. \\
\code{submit()} & attribue une séquence + push & Stabiliser l'ordre avant parallélisme. \\
\code{decodeFrame()} & valide protocole et CRC & Distinguer corruption de transport et trame valide. \\
\code{writerLoop()} & remet dans l'ordre & Garantir CSV/santé identiques avec 1 ou N workers. \\
\code{processApplicationPayload()} & détecte STGA & Garder la compatibilité avec les payloads opaques. \\
\code{TelemetryOutput::write()} & écrit CSV/JSON & Isoler les effets de bord disque du décodage parallèle. \\
\bottomrule
\end{longtable}

\begin{limitbox}{Arrêt}
Ctrl+C ou SIGTERM bascule \code{running=false}. Avec le timeout réseau par défaut, la boucle \code{poll()} réévalue l'arrêt au plus tard après environ 250 ms, hors temps de drainage en cours.
\end{limitbox}

\section{Commande UDP}
\begin{lstlisting}
./build/stgs_ground_station --udp --port 9000 --decoder-threads 4 --output telemetry.csv
\end{lstlisting}
Le début du pipeline est identique, mais \code{NetworkServer::run()} appelle \code{runUdp()}. Chaque datagramme est transmis comme un candidat complet. \code{MSG_TRUNC} permet de connaître la taille réelle : tout datagramme supérieur à 4123 octets est rejeté avant le codec.

\begin{callflow}{Différence essentielle UDP}
\ttfamily
recvfrom(MSG\_TRUNC)
$\rightarrow$ taille <= MaxFrameSize ?
$\rightarrow$ callback(ByteVector)
$\rightarrow$ pipeline commun
\end{callflow}

\source{src/NetworkServerUdp.cpp}{18--77}

\section{Sélection automatique d'un port local}
\begin{lstlisting}
./build/stgs_ground_station --tcp --auto-port 9000:9010 --output telemetry.csv
\end{lstlisting}
\code{parseArgs()} interdit la combinaison \code{--port} + \code{--auto-port}, exige TCP et limite la plage à 256 ports. Ensuite \code{run()} appelle \code{selectAutomaticPort()} avant de construire le logger et le pipeline.

\begin{callflow}{Auto-port station}
\ttfamily
selectAutomaticPort
$\rightarrow$ findFirstAvailableLocalPort
$\rightarrow$ checkLocalPortAvailability(port 9000)
$\rightarrow$ socket + bind temporaire
$\rightarrow$ si occupé : port suivant
$\rightarrow$ premier bind réussi = port choisi
$\rightarrow$ fermeture de la socket de test
$\rightarrow$ démarrage normal de la station sur ce port
\end{callflow}

\begin{warningbox}{Ce que \texttt{--auto-port} ne fait pas}
Il ne cherche pas une communication existante. Il cherche le premier port \textbf{libre pour bind} sur la machine locale. La recherche d'un listener existant est la fonction de \code{stgs_port_check} et de \code{--discover-ports} côté simulateur.
\end{warningbox}

\section{Replay STGF}
\begin{lstlisting}
./build/stgs_ground_station --replay frames.stgf --replay-rate 100 --output replay.csv
\end{lstlisting}
\code{parseArgs()} interdit le mélange replay + TCP/UDP. \code{receiveInput()} choisit \code{receiveReplay()}.

\begin{callflow}{Replay}
\ttfamily
FrameFileReader(path)
$\rightarrow$ validation magic STGF + version
$\rightarrow$ readNext()
$\rightarrow$ longueur u32 vérifiée [27,4123]
$\rightarrow$ lecture exacte de la trame
$\rightarrow$ submit()
$\rightarrow$ pipeline normal decode/CRC/santé/STGA/export
\end{callflow}

Avec \code{--replay-rate 0}, aucune temporisation n'est ajoutée. Au-dessus de 0, le code utilise \code{steady_clock} et \code{sleep_until} avec une échéance cumulée ; le temps du décodage n'est donc pas ajouté à chaque période.

\source{apps/ground_station/GroundStationInput.cpp}{30--63} ; \source{src/Replay.cpp}{103--147}

\section{Sortie CSV et JSON}
\subsection{CSV implicite}
\begin{lstlisting}
./build/stgs_ground_station --tcp --port 9000 --output telemetry.csv
\end{lstlisting}
L'extension autre que \code{.json} entraîne CSV en mode Auto.

\subsection{JSON implicite}
\begin{lstlisting}
./build/stgs_ground_station --tcp --port 9000 --output telemetry.json
\end{lstlisting}
L'extension exactement \code{.json} entraîne JSON.

\subsection{Forçage du format}
\begin{lstlisting}
./build/stgs_ground_station --tcp --port 9000 --output data.txt --output-format json
\end{lstlisting}
Le format demandé gagne sur l'extension.

\code{TelemetryOutput} ouvre le fichier en troncature dès la construction du pipeline. Chaque trame \textbf{validée} est exportée ; une trame rejetée CRC/format n'est pas écrite. Une erreur STGA interne ne retire pas la trame externe de l'export, car le framing et le CRC restent valides.

\source{apps/ground_station/TelemetryOutput.cpp}{23--72}

\section{Filtrage des signaux reçus}
\begin{lstlisting}
./build/stgs_ground_station --tcp --port 9000 --signal-filter none
./build/stgs_ground_station --tcp --port 9000 --signal-filter moving-average --filter-window 5
./build/stgs_ground_station --tcp --port 9000 --signal-filter sine-projection
\end{lstlisting}

\begin{tabularx}{\textwidth}{>{\bfseries}l X X}
\toprule
Mode & Fonction appelée & Conséquence \\
\midrule
none & copie des samples & Affiche le signal brut comme filtré. \\
moving-average & \code{movingAverageFilter()} & FIR centré, fenêtre impaire. \\
sine-projection & \code{sineProjectionFilter()} & Estime DC + coefficients sin/cos à la fréquence nominale. \\
\bottomrule
\end{tabularx}

Après filtrage, \code{computeSignalMetrics()} calcule RMS brut, RMS filtré, RMS du résidu, amplitude estimée et SNR de démonstration. Le résultat filtré est rendu dans le terminal ; le CSV/JSON conserve la trame reçue et son payload brut.

\section{Santé NOMINAL / DEGRADED}
Options principales :
\begin{lstlisting}
--degraded-window 100
--degraded-min-samples 20
--degraded-rejection-rate 0.10
--degraded-recovery-rate 0.03
--degraded-critical-count 8
--recovery-critical-count 3
\end{lstlisting}

Une trame CRC/format rejetée appelle \code{recordRejected()}. Une trame valide appelle \code{recordDecoded()}; CRITICAL et SAFE\_MODE comptent comme critiques. Aucune transition n'est évaluée avant \code{minSamples}.

\begin{callflow}{Hystérésis par défaut}
\ttfamily
NOMINAL
$\xrightarrow{rejets >= 10\% \; OU \; critiques >= 8}$ DEGRADED
$\xrightarrow{rejets <= 3\% \; ET \; critiques <= 3}$ NOMINAL
\end{callflow}

\source{src/StationHealth.cpp}{70--110}

\section{Logging, verbose et couleurs}
\begin{lstlisting}
./build/stgs_ground_station --tcp --port 9000 --log station.log --log-level info
./build/stgs_ground_station --tcp --port 9000 --verbose
./build/stgs_ground_station --tcp --port 9000 --no-color
\end{lstlisting}
\code{--verbose} active aussi le niveau Debug et journalise chaque trame décodée avec sa séquence. \code{--no-color} n'empêche pas la structure de l'UI ; il supprime seulement les codes ANSI.

\section{Table complète des bornes station}
\begin{longtable}{p{0.28\textwidth}p{0.2\textwidth}p{0.42\textwidth}}
\toprule
Option & Borne & Justification \\
\midrule
\endhead
\code{--port} & 1..65535 & Port 0 refusé à l'utilisateur. \\
\code{--auto-port} & <=256 ports & Diagnostic local volontairement borné. \\
\code{--decoder-threads} & 1..256 & Garde-fou CLI, pas limite protocole. \\
\code{--queue-capacity} & 1..1 000 000 & Évite zéro et saisies démesurées. \\
\code{--replay-rate} & >=0, fini & 0 = débit maximal. \\
\code{--filter-window} & impair, 1..1017 & Filtre centré ; borne dérivée du bloc signal maximal. \\
Fenêtre santé & 1..1 000 000 & Garde-fou d'allocation. \\
Probabilités santé & [0,1] & Taux mathématiquement valide. \\
\bottomrule
\end{longtable}

\end{document}
